\begin{exercise}{7}
Let $F\colon I\to\mathcal C$ be a functor. Show that a colimit of $F$
can equivalently be described as an initial cocone over $F$, and that
a limit of $F$ can be equivalently described as a terminal cone over $F$.

\begin{proof}
A cocone is an object $c$ with maps $\lambda_i\colon F(i)\to c$
satisfying $\lambda_jF(u)=\lambda_i$ for every $u\colon i\to j$.
A morphism $(c,\lambda)\to(c',\lambda')$ of cocones is a map
$h\colon c\to c'$ such that $h\lambda_i=\lambda'_i$ for all $i$.
These form the category of cocones over $F$.\footnote{The final
sentence of Observation~1.1.15 says that a colimit is initial in the
category of ``cones'', and a limit is terminal in the category of
``cocones''. The words must be interchanged: a colimit has maps
$F(i)\to c$ and hence is a cocone, whereas a limit has maps
$c\to F(i)$ and hence is a cone.}
Thus $(L,\lambda)$ is initial precisely when every compatible family
$\mu_i\colon F(i)\to c$ admits a unique $h\colon L\to c$ with
$h\lambda_i=\mu_i$. This is exactly the colimit property of
Definition~1.1.13.

Dually, a cone $(c,\pi)$ consists of compatible maps
$\pi_i\colon c\to F(i)$, and a morphism
$(c,\pi)\to(c',\pi')$ is a map $h\colon c\to c'$ with
$\pi'_ih=\pi_i$ for all $i$.
A cone $(L,p)$ is terminal precisely when every cone $(c,\pi)$
admits a unique $h\colon c\to L$ satisfying $p_ih=\pi_i$.
This is exactly the limit property.
\end{proof}
\end{exercise}
